\documentclass[10pt,a4paper]{amsart}
\usepackage{amsmath,amssymb,amsthm,graphicx}
% Retained preamble for the Stacks Project chapter morphisms.tex (Morphisms of Schemes).
% The chapter file pinned-original-source.tex carries only \input{preamble} on its first line, so no
% preamble text lives inside the chapter itself.  The Stacks Project uses one single shared preamble for
% every chapter; the snapshot stages it verbatim in sources/stacks-chow/preamble.tex.  The block below is
% copied mechanically, line for line, from that file's lines 14-18 (the Xy-pic setup) and lines 154-217
% (the native theorem declarations together with the mathematical macros).  Nothing was added, removed or
% redefined.  The wrapper class and the amsmath/amssymb/amsthm/graphicx base packages are supplied by the
% assembler, so the class/\externaldocument part of the shared preamble is not repeated here.
% For commutative diagrams we use Xy-pic
\usepackage[all]{xy}

% We use 2cell for 2-commutative diagrams.
\xyoption{2cell}
% Theorem environments.
%
\theoremstyle{plain}
\newtheorem{theorem}[subsection]{Theorem}
\newtheorem{proposition}[subsection]{Proposition}
\newtheorem{lemma}[subsection]{Lemma}

\theoremstyle{definition}
\newtheorem{definition}[subsection]{Definition}
\newtheorem{example}[subsection]{Example}
\newtheorem{exercise}[subsection]{Exercise}
\newtheorem{situation}[subsection]{Situation}

\theoremstyle{remark}
\newtheorem{remark}[subsection]{Remark}
\newtheorem{remarks}[subsection]{Remarks}

\numberwithin{equation}{subsection}

% Macros
%
\def\lim{\mathop{\mathrm{lim}}\nolimits}
\def\colim{\mathop{\mathrm{colim}}\nolimits}
\def\Spec{\mathop{\mathrm{Spec}}}
\def\Hom{\mathop{\mathrm{Hom}}\nolimits}
\def\Ext{\mathop{\mathrm{Ext}}\nolimits}
\def\SheafHom{\mathop{\mathcal{H}\!\mathit{om}}\nolimits}
\def\SheafExt{\mathop{\mathcal{E}\!\mathit{xt}}\nolimits}
\def\Sch{\mathit{Sch}}
\def\Mor{\mathop{\mathrm{Mor}}\nolimits}
\def\Ob{\mathop{\mathrm{Ob}}\nolimits}
\def\Sh{\mathop{\mathit{Sh}}\nolimits}
\def\NL{\mathop{N\!L}\nolimits}
\def\CH{\mathop{\mathrm{CH}}\nolimits}
\def\proetale{{pro\text{-}\acute{e}tale}}
\def\etale{{\acute{e}tale}}
\def\QCoh{\mathit{QCoh}}
\def\Ker{\mathop{\mathrm{Ker}}}
\def\Im{\mathop{\mathrm{Im}}}
\def\Coker{\mathop{\mathrm{Coker}}}
\def\Coim{\mathop{\mathrm{Coim}}}

% Boxtimes
%
\DeclareMathSymbol{\boxtimes}{\mathbin}{AMSa}{"02}

%
% Macros for moduli stacks/spaces
%
\def\QCohstack{\mathcal{QC}\!\mathit{oh}}
\def\Cohstack{\mathcal{C}\!\mathit{oh}}
\def\Spacesstack{\mathcal{S}\!\mathit{paces}}
\def\Quotfunctor{\mathrm{Quot}}
\def\Hilbfunctor{\mathrm{Hilb}}
\def\Curvesstack{\mathcal{C}\!\mathit{urves}}
\def\Polarizedstack{\mathcal{P}\!\mathit{olarized}}
\def\Complexesstack{\mathcal{C}\!\mathit{omplexes}}
% \Pic is the operator that assigns to X its picard group, usage \Pic(X)
% \Picardstack_{X/B} denotes the Picard stack of X over B
% \Picardfunctor_{X/B} denotes the Picard functor of X over B
\def\Pic{\mathop{\mathrm{Pic}}\nolimits}
\def\Picardstack{\mathcal{P}\!\mathit{ic}}
\def\Picardfunctor{\mathrm{Pic}}
\def\Deformationcategory{\mathcal{D}\!\mathit{ef}}

% --- appended 2026-10-12 (Chief): the project own definitions of the reference/slogan/history
% display environments, copied mechanically from the shared Stacks preamble lines 8-12
% (\usepackage{verbatim} plus the three \newenvironment lines, which the chapter preamble keeps
% verbatim from the project preamble). Nothing added, removed or redefined: these environments
% carry no mathematical content and the project itself typesets them as comments. Without these
% three lines a counted statement whose own text uses them fails with Environment slogan undefined.
\usepackage{verbatim}
\newenvironment{reference}{\comment}{\endcomment}
\newenvironment{slogan}{\comment}{\endcomment}
\newenvironment{history}{\comment}{\endcomment}
\usepackage{cleveref}
\crefname{originallemma}{original lemma}{original lemmas}
\Crefname{originallemma}{Original lemma}{Original lemmas}
\makeatletter
\newlabel{schemes-lemma-points-fibre-product}{{01JT}{source}{}{sourcepointer-1}{}}
\newlabel{schemes-lemma-points-fibre-product@cref}{{[originallemma][0][]01JT}{[1][0][]source}{}{}{}}
\makeatother
\crefname{originallemma}{original lemma}{original lemmas}
\Crefname{originallemma}{Original lemma}{Original lemmas}
\makeatletter
\newlabel{schemes-lemma-quasi-compact-affine}{{01K4}{source}{}{sourcepointer-2}{}}
\newlabel{schemes-lemma-quasi-compact-affine@cref}{{[originallemma][0][]01K4}{[1][0][]source}{}{}{}}
\makeatother
\title{A universally closed morphism of schemes is quasi-compact}
\author{The Stacks Project Authors}
\date{stacks-project:morphisms.tex:local-snapshot-20261012; GNU FDL 1.2 or later}
\begin{document}
\maketitle
\noindent\textit{Editorial note.} All delivered mathematical text is the author's own, copied from the pinned Stacks Project chapter file pinned-original-source.tex; no statement or proof step was rewritten, reordered or removed.\par\medskip
\noindent\textit{Editorial note.} Following the multi-claim convention for this source, source\_url is the source-level URL https://stacks.math.columbia.edu/ and the source label lemma-universally-closed-quasi-compact is carried in claim\_key and here rather than in a per-claim tag URL.\par\medskip
\noindent\textit{Editorial note.} The preamble is the chapter's shared Stacks Project preamble, retained as sources/stacks-morphisms/retained-preamble.tex: a line-for-line excerpt of the shared Stacks preamble staged with this snapshot. No macro was invented or redefined, and this card adds no preamble text of its own.\par\medskip
\noindent\textit{Editorial note.} Statement span 9985-9991 and proof span 9993-10047 come from the halfA pool card operations/pool-cards/stacks-morphisms-halfA.json and were re-read in the pinned chapter file before this card was written.\par\medskip
\noindent\textit{Editorial note.} The statement contains a reference environment attributing the result to Bjorn Poonen; the project's own shared preamble (retained in full at 88 lines) defines that environment as a comment environment, so the attribution is not typeset, exactly as in the project's own output, and no bibliography is attached.\par\medskip
\noindent\textit{Editorial note.} The proof names an auxiliary affine space with one coordinate per affine open of the source and writes the corresponding closed subset explicitly; the author's notation is preserved unchanged.\par\medskip
\noindent\textit{Editorial note.} 2 cross-references occur inside the counted spans; they are listed in external\_references as navigation pointers only.\par\medskip
\noindent\textit{Allowed standard prerequisites (editorial).}\begin{itemize}
\item The characterization of quasi-compactness by the existence of an affine open whose preimage is not quasi-compact, cited from the schemes chapter. Cited at proof lines 9995-9997 and again at 10011-10013. Produces the affine base to work over and the bad point at which every neighbourhood has non-quasi-compact preimage.
\item The description of points of a fibre product by a pair of points together with a prime in the tensor product of their residue fields, cited from the schemes chapter. Cited at proof lines 10038-10040. Allows the point constructed over a chosen point of X and a chosen point of the auxiliary affine space to be lifted to the fibre product.
\item The description of the closed subscheme cut out by a family of coordinate functions as a polynomial algebra on the remaining coordinates, and the bijection between its points over a given base point and the spectrum of the same polynomial algebra over that residue field. Used at proof lines 10033-10037. Makes it possible to choose the auxiliary point avoiding the zero locus of a prescribed polynomial while lying in the vanishing loci of the omitted coordinates.
\end{itemize}\medskip
\ifdefined\hypertarget\hypertarget{sourcepointer-1}{}\fi
\noindent\textit{Source reference.} Original lemma 01JT: focused-seeds/AGP/sources/stacks-snapshot-20261011/tags.txt, source lines 2026--2026. Same-chapter reference used in the retained author text at source line(s) 1326, 1485, 1652, 3643, 10040. It is supplied only as a navigation pointer so that the retained text compiles; the counted item delivers the author's own statement and the author's complete proof, and no part of the referenced item is copied in or substituted for the counted argument.\par\medskip
\ifdefined\hypertarget\hypertarget{sourcepointer-2}{}\fi
\noindent\textit{Source reference.} Original lemma 01K4: focused-seeds/AGP/sources/stacks-snapshot-20261011/tags.txt, source lines 2037--2037. Same-chapter reference used in the retained author text at source line(s) 1727, 1761, 2213, 2252, 9996, 10014. It is supplied only as a navigation pointer so that the retained text compiles; the counted item delivers the author's own statement and the author's complete proof, and no part of the referenced item is copied in or substituted for the counted argument.\par\medskip
\begin{lemma}
\label{lemma-universally-closed-quasi-compact}
\begin{reference}
Due to Bjorn Poonen.
\end{reference}
A universally closed morphism of schemes is quasi-compact.
\end{lemma}

\begin{proof}
Let $f : X \to S$ be a morphism. Assume that $f$ is not quasi-compact.
Our goal is to show that $f$ is not universally closed. By
Schemes, Lemma \ref{schemes-lemma-quasi-compact-affine}
there exists an affine open $V \subset S$ such that $f^{-1}(V)$ is
not quasi-compact. To achieve our goal it suffices to show that
$f^{-1}(V) \to V$ is not universally closed, hence we may assume that
$S = \Spec(A)$ for some ring $A$.

\medskip\noindent
Write $X = \bigcup_{i \in I} X_i$ where the $X_i$ are affine open subschemes
of $X$.  Let $T = \Spec(A[y_i ; i \in I])$.
Let $T_i = D(y_i) \subset T$. Let $Z$ be the closed set
$(X \times_S T) - \bigcup_{i \in I} (X_i \times_S T_i)$.  It suffices to
prove that the image $f_T(Z)$ of $Z$ under $f_T : X \times_S T \to T$
is not closed.

\medskip\noindent
There exists a point $s \in S$ such that there is no
neighborhood $U$ of $s$ in $S$ such that $X_U$ is quasi-compact.
Otherwise we could cover $S$ with finitely many such $U$ and
Schemes, Lemma \ref{schemes-lemma-quasi-compact-affine}
would imply $f$ quasi-compact. Fix such an $s \in S$.

\medskip\noindent
First we check that $f_T(Z_s) \ne T_s$.  Let $t \in T$ be the point
lying over $s$ with $\kappa(t) = \kappa(s)$ such that $y_i = 1$ in
$\kappa(t)$ for all $i$.
Then $t \in T_i$ for all $i$, and the fiber of $Z_s \to T_s$ above
$t$ is isomorphic to $(X - \bigcup_{i \in I} X_i)_s$, which is empty.
Thus $t \in T_s - f_T(Z_s)$.

\medskip\noindent
Assume $f_T(Z)$ is closed in $T$. Then there exists an element
$g \in A[y_i; i \in I]$ with $f_T(Z) \subset V(g)$ but $t \not \in V(g)$.
Hence the image of $g$ in $\kappa(t)$ is nonzero. In particular some
coefficient of $g$ has nonzero image in $\kappa(s)$. Hence this coefficient is
invertible on some neighborhood $U$ of $s$. Let $J$ be the finite set of
$j \in I$ such that $y_j$ appears in $g$. Since $X_U$ is not quasi-compact,
we may choose a point $x \in X - \bigcup_{j \in J} X_j$ lying above some
$u \in U$. Since $g$ has a coefficient that is invertible on $U$, we can
find a point $t' \in T$ lying above $u$ such that $t' \not \in V(g)$ and
$t' \in V(y_i)$ for all $i \notin J$. This is true because
$V(y_i; i \in I, i \not\in J) = \Spec(A[y_j; j\in J])$
and the set of points of this scheme lying over $u$ is bijective
with $\Spec(\kappa(u)[y_j; j \in J])$. In other words $t' \notin T_i$
for each $i \notin J$. By
Schemes, Lemma \ref{schemes-lemma-points-fibre-product}
we can find a point $z$ of $X \times_S T$ mapping to $x \in X$ and to
$t' \in T$. Since $x \not \in X_j$ for $j \in J$ and $t' \not \in T_i$
for $i \in I \setminus J$ we see that $z \in Z$. On the other hand
$f_T(z) = t' \not \in V(g)$ which contradicts $f_T(Z) \subset V(g)$.
Thus the assumption ``$f_T(Z)$ closed'' is wrong and we conclude indeed
that $f_T$ is not closed, as desired.
\end{proof}

\end{document}
